\documentclass[11pt]{article}
\usepackage[a4paper,margin=26mm]{geometry}
\usepackage{amsmath,amssymb,amsthm}
\usepackage[T1]{fontenc}
\usepackage{booktabs}
\usepackage{array}
\usepackage{alltt}
\usepackage{xurl}
\usepackage{hyperref}
\hypersetup{colorlinks=true,linkcolor=blue,urlcolor=blue,citecolor=blue,
 pdftitle={Complete mappings of the Brauer and partial Brauer monoids},
 pdfauthor={Horace Dong}}
\newtheorem{theorem}{Theorem}[section]
\newtheorem{lemma}[theorem]{Lemma}
\newtheorem{proposition}[theorem]{Proposition}
\newtheorem{corollary}[theorem]{Corollary}
\theoremstyle{remark}
\newtheorem{remark}[theorem]{Remark}
\newcommand{\B}{\mathcal B}
\newcommand{\PB}{\mathcal{PB}}
\newcommand{\M}{\mathcal M}
\newcommand{\Z}{\mathbb Z}
\newcommand{\tp}{\operatorname{top}}
\newcommand{\bt}{\operatorname{bot}}
\newcommand{\rk}{\operatorname{rank}}
\newcommand{\Mt}{\M_2}
\title{Complete mappings of the Brauer\\ and partial Brauer monoids}
\author{Horace Dong\\ \small Independent researcher\\
\small\texttt{maxwelldhx@gmail.com}\\
\small ORCID: \href{https://orcid.org/0009-0002-7554-7548}{0009-0002-7554-7548}}
\date{September 2026}
\begin{document}
\maketitle
\begin{abstract}
A complete mapping of a semigroup $S$ is a bijection $\alpha\colon S\to S$
such that $x\mapsto x\cdot x\alpha$ is also a bijection. Problem~15.11 of
Ara\'ujo, Bentz, Cameron, Hendrey and Kinyon (arXiv:2608.25092) asks for a
classification of complete mappings of the Brauer monoid $\B_n$ and the
partial Brauer monoid $\PB_n$. By their reduction to principal factors and the
Hall--Paige theorem, the open cases are the principal factors of ranks $2$
and~$3$, together with the question of which of the monoids have a complete
mapping. We settle the existence question in all cases. Every proper
principal factor of $\B_n$ has a complete mapping. Every proper principal
factor of $\PB_n$ has one, except the rank-$2$ factor of $\PB_3$, which is the
Brandt semigroup $B(S_2,3)$ and has none. Consequently $\B_n$ and $\PB_n$ have
a complete mapping if and only if $n\notin\{2,3\}$, as for the full
transformation monoid and the partition monoid. The new case is that of $\B_n$
with $n\equiv2,3\pmod 4$, in which the rank-$2$ or rank-$3$ factor has an odd
number of $\mathcal L$-classes. For it we give an explicit complete mapping,
built from a triangle of half-diagrams whose sandwich permutations have odd
product and from a perfect matching of the remaining half-diagrams. The other
cases follow directly from results of Ara\'ujo et al.\ once the Rees structure
of the principal factors is written down. The two explicit constructions of
complete mappings of Rees $0$-matrix semigroups over $\Z_2$ on which the
proofs rest are formally verified in Lean~4, for arbitrary finite index sets;
the rest of the argument is not formalised. We also give explicit complete
mappings of $\B_n$ and $\PB_n$ for $n\in\{1,4,5,6\}$ and of the two smallest
factors in the new case, each checked by a program that shares no code with
the program that built it. The proofs were found by AI coding agents working
under the author's direction; their role is described at the end of the
paper.
\end{abstract}

\section{Introduction}

\paragraph{Background.}
Let $(S,\cdot)$ be a magma. A \emph{complete mapping} of $S$ is a bijection
$\alpha\colon S\to S$ such that the map $\theta\colon S\to S$,
$x\theta=x\cdot x\alpha$, is also a bijection. As in~\cite{abchk}, maps act on
the right. For a finite magma this is the same as a transversal of the Cayley
table. For finite groups, existence is decided by the Hall--Paige theorem, in
the form recorded as \cite[Theorem~1.1]{abchk}: a finite group has a complete
mapping if and only if its Sylow $2$-subgroups are trivial or non-cyclic. In
particular, the symmetric group $S_r$ has a complete mapping if and only if
$r\notin\{2,3\}$.

Ara\'ujo, Bentz, Cameron, Hendrey and Kinyon~\cite{abchk} developed the theory
for finite semigroups. A finite semigroup has a complete mapping if and only
if each of its principal factors has one \cite[Theorem~4.4]{abchk}, which
reduces the question to Rees $0$-matrix semigroups; these are studied in
detail in \cite[Sections~6 and~7]{abchk}. Among the applications, the full
transformation monoid $T_n$ and the partition monoid $\mathcal P_n$ have
complete mappings if and only if $n=1$ or $n\ge4$
\cite[Theorems~9.16 and~8.5]{abchk}, the symmetric inverse monoid $I_n$ has one
if and only if $n\equiv0,1\pmod4$ \cite[Corollary~10.2]{abchk}, and the planar
partition, Motzkin and Jones monoids have complete mappings
\cite[Corollary~8.8]{abchk}. The Brauer monoid $\B_n$ and the partial Brauer
monoid $\PB_n$ were left open. Problem~15.11 of~\cite{abchk} asks for a
classification of their complete mappings; it notes that, by the Hall--Paige
theorem, the proper principal factors still to be decided are those of ranks
$2$ and~$3$, and it also asks for which $n$ the whole monoid has a complete
mapping. Here a principal factor is \emph{proper} if it is not the principal
factor of the group of units.

\paragraph{Results.}
We decide every case.

\begin{theorem}\label{thm:main}
Let $n\ge1$.
\begin{itemize}
\item[(a)] Every proper principal factor of $\B_n$ has a complete mapping.
\item[(b)] Every proper principal factor of $\PB_n$ has a complete mapping,
 except the principal factor of rank~$2$ of $\PB_3$. That factor is isomorphic
 to the Brandt semigroup $B(S_2,3)=\M^0(S_2,[3],[3],\Delta_3)$, where
 $\Delta_3$ is the $3\times3$ identity matrix, and it has no complete mapping.
\item[(c)] The following are equivalent: (i) $\B_n$ has a complete mapping;
 (ii) $\PB_n$ has a complete mapping; (iii) $S_n$ has a complete mapping;
 (iv) $n=1$ or $n\ge4$.
\end{itemize}
\end{theorem}

The principal factor of rank $r$ is a Rees $0$-matrix semigroup
$\M^0(S_r,H_r,H_r,P)$ over the symmetric group $S_r$, indexed on both sides by
the set $H_r$ of half-diagrams of rank~$r$ (Proposition~\ref{prop:rees}). Its
sandwich matrix has identity diagonal. So when $S_r$ has a complete mapping,
that is, when $r\notin\{2,3\}$, the factor has one by
\cite[Theorem~6.2]{abchk}. When $r\in\{2,3\}$ and $N=|H_r|$ is even, it has one
by \cite[Corollary~7.6]{abchk}. The remaining factors are those with
$r\in\{2,3\}$ and $N$ odd, the case in which both index sets are odd and, as
\cite[Section~7]{abchk} explains, the images of the sandwich entries in a
cyclic $2$-quotient of $S_r$ matter, not just the pattern of non-zero entries.
Table~\ref{tab:cases} lists the parities. For $\PB_n$ the only proper factor
with $r\in\{2,3\}$ and $N$ odd is the rank-$2$ factor of $\PB_3$, and it has no
complete mapping by \cite[Theorem~7.19]{abchk}. For $\B_n$ there are
infinitely many such factors, those of rank~$2$ for $n\equiv2\pmod4$ and of
rank~$3$ for $n\equiv3\pmod4$; the first are the rank-$2$ factor of $\B_6$,
with $N=45$ and $4050$ non-zero elements, and the rank-$3$ factor of $\B_7$,
with $N=105$ and $66150$ non-zero elements. For these factors we construct
complete mappings explicitly (Section~\ref{sec:odd}), from three pairwise
compatible half-diagrams whose sandwich permutations have odd product and a
perfect matching of the other half-diagrams by compatible pairs.

\begin{table}[ht]\centering\small
\begin{tabular}{@{}l l l l >{\raggedright\arraybackslash}p{0.33\textwidth}@{}}
\toprule
Monoid & rank $r$ & $N=|H_r|$ & $N$ odd exactly when & complete mapping of the factor\\
\midrule
$\B_n$ & $2$ & $\binom n2\,(n-3)!!$ & $n\equiv2\pmod4$ &
 yes: \cite[Corollary~7.6]{abchk} if $N$ is even, Section~\ref{sec:odd} if
 $N$ is odd\\\addlinespace
$\B_n$ & $3$ & $\binom n3\,(n-4)!!$ & $n\equiv3\pmod4$ &
 yes: as for rank $2$\\\addlinespace
$\PB_n$ & $2,3$ & $\binom nr\,a(n-r)$ & $(n,r)=(3,2)$ &
 yes if $N$ is even \cite[Corollary~7.6]{abchk}; none for $(n,r)=(3,2)$
 (Proposition~\ref{prop:pb3})\\
\bottomrule
\end{tabular}
\caption{The proper principal factors of ranks $2$ and $3$ ($r<n$, and
$r\equiv n\pmod 2$ for $\B_n$). Here $N$ is the number of $\mathcal
L$-classes of the factor, which equals the number of $\mathcal R$-classes, and
$a(k)$ is the number of involutions of a $k$-set. The factors of rank
$r\notin\{2,3\}$ have complete mappings by \cite[Theorem~6.2]{abchk}.}
\label{tab:cases}
\end{table}

All our complete mappings are given by explicit formulas
(Lemmas~\ref{lem:A}--\ref{lem:C}). The two constructions for ranks $2$ and $3$
(Lemmas~\ref{lem:B} and~\ref{lem:C}), stated for a Rees $0$-matrix semigroup
over $\Z_2$ with identity diagonal and an arbitrary finite index set, are
formally verified in Lean~4 (Section~\ref{sec:lean}). The facts about diagrams
(Sections~\ref{sec:rees}, \ref{sec:odd} and~\ref{sec:pb}) and the results of
\cite{abchk} that we use are not formalised; for them the evidence is the
written proofs, supported by the computer checks of Section~\ref{sec:comp}.

\paragraph{Scope.}
We read ``classify complete mappings'' in the sense explained by the second
sentence of Problem~15.11: decide, for every $n$ and every rank~$r$, whether the
principal factor of rank~$r$ has a complete mapping, and decide whether the
whole monoid has one. Theorem~\ref{thm:main} does this completely. We do not
describe the set of all complete mappings, as \cite[Proposition~6.5]{abchk}
does for the combinatorial Brandt semigroups by means of Latin squares.

\paragraph{Prior work.}
Problem~15.11 appears in version~1 of~\cite{abchk}, submitted to the arXiv on
25 August 2026, which was still the only version when we checked on
26~September 2026. We found no solution of, and no announced progress on,
Problem~15.11 in the following sources, which we searched on 26~September 2026
between 06:49 and 07:46 UTC: the arXiv (the listing of~\cite{abchk} and eight
title and abstract searches through the arXiv API); GitHub searches of issues,
pull requests, commits, code and repositories; the public repositories of
AI-assisted work on open problems that we check, namely the file tree, issues
and pull requests of the \texttt{trureturing} repository~\cite{trurepo}, and
the repositories SCOPE2026, AI-Has-Taste, formal-conjectures and astrafala; the
problem database MathDB (\url{https://mathdb.com}); MathOverflow and Mathematics
Stack Exchange, through their search interface; Crossref; web search; the blog
post announcing~\cite{abchk} and its comments~\cite{camblog}, and the
publication list and arXiv listing of the third author; and the slides of a
2025 talk on the subject by the fifth author~\cite{kinyon}. The only earlier
AI-assisted work on~\cite{abchk} that we found is issue~\#9377 and pull
request~\#9405 of~\cite{trurepo}, which concern Problem~15.5 and do not mention
Problem~15.11 or diagram monoids. Our searches on 26~September 2026 for works
citing~\cite{abchk} in OpenAlex and Semantic Scholar failed (HTTP error 429,
rate limit exceeded), so citing works were not checked systematically; we did
not search Google Scholar, zbMATH or MathSciNet. A final check on 26 September
2026 between 15:49 and 15:52 UTC (the arXiv, where~\cite{abchk} still had only
version~1; GitHub; MathDB; and MathOverflow) found nothing new. This reports
the coverage of our searches, not a guarantee of priority. Our results are a direct extension
of the methods of~\cite{abchk}: the reduction to principal factors, the
criteria for Rees $0$-matrix semigroups, the lifting lemma and the
non-existence theorem that we use are theirs, and our treatment of $\B_n$
follows their proof for the partition monoid \cite[Theorem~8.5]{abchk},
including the triangle of half-diagrams in Lemma~\ref{lem:triangle}.

\paragraph{Organisation.}
Section~\ref{sec:prelim} fixes notation and lists the results
of~\cite{abchk} that we use. Section~\ref{sec:rees} gives the Rees structure of
the principal factors of $\B_n$ and $\PB_n$. Section~\ref{sec:cm} gives
explicit complete mappings of Rees $0$-matrix semigroups with identity
diagonal. Section~\ref{sec:odd} treats $\B_n$ and proves
Theorem~\ref{thm:main}(a), Section~\ref{sec:pb} treats $\PB_n$ and proves
Theorem~\ref{thm:main}(b), and Section~\ref{sec:whole} proves
Theorem~\ref{thm:main}(c). Section~\ref{sec:comp} reports the computations and
Section~\ref{sec:lean} the formal verification.

\section{Preliminaries}\label{sec:prelim}

\paragraph{Rees $0$-matrix semigroups.}
We follow \cite[Section~4]{abchk}. Let $G$ be a finite group, let $I$ and
$\Lambda$ be non-empty finite sets, and let $P=(p_{\lambda i})$ be a
$\Lambda\times I$ matrix over $G\cup\{0\}$ with no zero row or column. The
Rees $0$-matrix semigroup $\M^0(G,I,\Lambda,P)$ is $(I\times G\times\Lambda)
\cup\{0\}$ with the product
\[
 (i,a,\lambda)(j,b,\mu)=\begin{cases}(i,\,a\,p_{\lambda j}\,b,\,\mu)&
 \text{if }p_{\lambda j}\ne0,\\ 0&\text{otherwise,}\end{cases}
\]
and $0x=x0=0$. The \emph{pattern} of $P$ is the zero-one matrix obtained by
replacing every element of $G$ by~$1$. For a $\mathcal J$-class $J$ of a finite
semigroup $S$, the principal factor $J^0$ is $J\cup\{0\}$, with the product of
$S$ when it lies in $J$ and $0$ otherwise.

\paragraph{Results of~\cite{abchk} that we use.}
We use the following results, with the numbering of~\cite{abchk}.
\begin{itemize}
\item[(R1)] \cite[Theorem~1.1]{abchk} (Hall--Paige). A finite group has a
 complete mapping if and only if its Sylow $2$-subgroups are trivial or
 non-cyclic.
\item[(R2)] \cite[Proposition~1.4]{abchk}. A complete mapping of a magma with
 zero fixes the zero, and so does the corresponding map $\theta$.
\item[(R3)] \cite[Corollary~3.4]{abchk}. If a finite monoid has a complete
 mapping, then so does its group of units.
\item[(R4)] \cite[Theorem~4.4]{abchk}. A finite semigroup has a complete mapping
 if and only if every principal factor $J^0$ has one.
\item[(R5)] \cite[Lemma~6.3]{abchk}, which extends the lifting construction of
 \cite[Theorem~5.2]{abchk} to Rees $0$-matrix semigroups. Let $N$ be a normal
 subgroup of $G$ that has a complete mapping, and let $\bar P$ be the entrywise
 image of $P$ in $G/N$, zero entries staying zero. If $\M^0(G/N,I,\Lambda,\bar
 P)$ has a complete mapping, then so does $\M^0(G,I,\Lambda,P)$.
\item[(R6)] \cite[Theorem~6.2]{abchk}, implication (c)$\Rightarrow$(a). If $G$
 has a complete mapping and, for every $s$ rows of the pattern, at least
 $s|I|/|\Lambda|$ columns have a non-zero entry in one of these rows, then
 $\M^0(G,I,\Lambda,P)$ has a complete mapping. If $|I|=|\Lambda|$ and the pattern
 contains a one-transversal (a permutation matrix), the condition holds.
\item[(R7)] \cite[Corollary~7.6]{abchk}. If $G$ has no complete mapping,
 $|I|=|\Lambda|$ is even and the pattern of $P$ contains a one-transversal,
 then $\M^0(G,I,\Lambda,P)$ has a complete mapping.
\item[(R8)] \cite[Theorem~7.19]{abchk}. If $G$ has non-trivial cyclic Sylow
 $2$-subgroups, $|I|$ and $|\Lambda|$ are odd, and $P$ is obtained from a
 normalised matrix over $G$ in which every entry has odd order by replacing
 some entries by~$0$, then $\M^0(G,I,\Lambda,P)$ has no complete mapping.
\end{itemize}
By (R1), $S_r$ has a complete mapping for $r\le1$ (the trivial group) and for
$r\ge4$ (its Sylow $2$-subgroups contain $\langle(1\,2),(3\,4)\rangle\cong
C_2\times C_2$), and none for $r\in\{2,3\}$ (the Sylow $2$-subgroups have
order~$2$).

\paragraph{Diagrams.}
Let $[n]=\{1,\dots,n\}$ and $[n]'=\{1',\dots,n'\}$. The partial Brauer monoid
$\PB_n$ consists of the partitions of $[n]\cup[n]'$ into blocks of size $1$
or~$2$, and the Brauer monoid $\B_n$ of those with all blocks of size~$2$; see
\cite[Section~2.2]{emrt}. We draw $[n]$ as the top row and $[n]'$ as the
bottom row. The product $xy$ is formed by drawing $x$ above~$y$, identifying
the bottom row of $x$ with the top row of $y$ (the \emph{middle row}), keeping
the connectivity of the outer points, and discarding closed loops and
components contained in the middle row. A block meeting both rows is a
\emph{through block}; the \emph{rank} of $x$ is its number of through blocks,
and $\rk(xy)\le\min(\rk x,\rk y)$. As recalled in \cite[Section~15]{abchk},
it is proved in \cite[Proposition~2.1 and Remark~2.2(v)--(vii)]{emrt} that the
$\mathcal J$-classes of $\B_n$ and $\PB_n$ are the rank classes and that a
maximal subgroup in rank~$r$ is isomorphic to~$S_r$. We write $J_r$ for the
set of elements of rank~$r$; in $\B_n$ it is non-empty exactly when $r\le n$
and $r\equiv n\pmod2$. The group of units is $J_n$, the set of permutation
diagrams, isomorphic to $S_n$.

\section{Rees coordinates}\label{sec:rees}
Let $D\in\{\B_n,\PB_n\}$ and let $r$ be a rank that occurs in~$D$. A
\emph{half-diagram} of rank~$r$ is a pair $h=(F,M)$, where $F=F(h)\subseteq[n]$
has $|F|=r$ (the \emph{free points}) and $M=M(h)$ is a set of disjoint
$2$-subsets (\emph{arcs}) of $[n]\setminus F$ covering $[n]\setminus F$ if
$D=\B_n$, and arbitrary if $D=\PB_n$; the points of $[n]\setminus F$ not
covered by $M$ are \emph{singletons}. Let $H=H_r(D)$ be the set of
half-diagrams of rank~$r$ and $N=|H|$. Then
\[
 N=\binom nr(n-r-1)!!\ \text{ for }\B_n,\qquad N=\binom nr\,a(n-r)\ \text{ for
 }\PB_n,
\]
where $(-1)!!=1$ and $a(k)$ is the number of involutions of a $k$-set.

For $x\in J_r$, let $\tp(x)=(F,M)$, where $F$ is the set of top points of
through blocks and $M$ the set of blocks contained in the top row; let
$\bt(x)$ be defined in the same way from the bottom row, with primes removed.
If $a_1<\dots<a_r$ are the top points and $b_1<\dots<b_r$ the bottom points
(without primes) of the through blocks, define $g(x)\in S_r$ by requiring that
$\{a_k,b'_{k\,g(x)}\}$ is a block of~$x$ for every~$k$. Permutations act on the
right, and $gh$ means first $g$, then~$h$.

For $\lambda,i\in H$, let $\Gamma(\lambda,i)$ be the multigraph on $[n]$ whose
edges are the arcs of $\lambda$ and the arcs of~$i$ (an arc of both gives a
double edge). Every vertex has degree at most~$2$, so the components are paths
and cycles. A point of $F(\lambda)$ lies on no arc of $\lambda$, so it has
degree at most~$1$ and is an end of a path. We say that $\lambda$ and $i$ are
\emph{compatible}, and write $\lambda\sim i$, if the path starting at each
point of $F(\lambda)$ ends at a point of $F(i)$; a point of $F(\lambda)\cap
F(i)$ is a path of length~$0$. The $r$ paths are distinct components, so they
then match $F(\lambda)$ bijectively with $F(i)$, and we define $p_{\lambda
i}\in S_r$ by requiring that the $k$-th smallest point of $F(\lambda)$ is joined
to the $(k\,p_{\lambda i})$-th smallest point of~$F(i)$.

\begin{proposition}\label{prop:rees}
Let $D\in\{\B_n,\PB_n\}$, let $r$ be a rank that occurs in $D$, and let
$H=H_r(D)$.
\begin{itemize}
\item[(a)] The map $x\mapsto(\tp(x),g(x),\bt(x))$ is a bijection from $J_r$ to
 $H\times S_r\times H$.
\item[(b)] For $x,y\in J_r$ we have $\rk(xy)=r$ if and only if
 $\bt(x)\sim\tp(y)$, and then $\tp(xy)=\tp(x)$, $\bt(xy)=\bt(y)$ and
 $g(xy)=g(x)\,p_{\bt(x),\tp(y)}\,g(y)$.
\item[(c)] For all $\lambda,i\in H$: $\lambda\sim\lambda$ and
 $p_{\lambda\lambda}=1$; and $\lambda\sim i$ if and only if $i\sim\lambda$, in
 which case $p_{i\lambda}=p_{\lambda i}^{-1}$.
\item[(d)] If $F(\lambda)=F(i)$, then $\lambda\sim i$ and $p_{\lambda i}=1$.
\end{itemize}
Consequently $J_r^0\cong\M^0(S_r,H,H,P)$, where $P$ is the $H\times H$ matrix
with entries $p_{\lambda i}$ if $\lambda\sim i$ and $0$ otherwise; the
isomorphism extends the map in~(a) by $0\mapsto0$.
\end{proposition}

\begin{proof}
(a) An element of rank $r$ is determined by its blocks inside the top row, its
blocks inside the bottom row, and a bijection from its top through points to
its bottom through points, and every such choice occurs. The first two data
are $\tp(x)$ and $\bt(x)$, and $g(x)$ encodes the bijection.

(b) The blocks of $x$ inside the top row are blocks of~$xy$. Consider the top
point $a_k$ of a through block of $x$. Its strand reaches the middle row at
$b=b_{k\,g(x)}\in F(\bt(x))$ and continues along the path of
$\Gamma(\bt(x),\tp(y))$ starting at~$b$. This path alternates between arcs of
$\tp(y)$ and arcs of $\bt(x)$, beginning with an arc of $\tp(y)$, and stops at
a point $e$ at which the next arc is missing. So either $e\in F(\tp(y))$, and
the strand continues through $y$ to the bottom row; or $e\ne b$ and $e\in
F(\bt(x))$, and $a_k$ is joined to another top through point of~$x$; or $e$ is
a singleton of $\bt(x)$ or of $\tp(y)$ (only in $\PB_n$), and $a_k$ becomes a
singleton. Every through block of $xy$ contains a top point of a through block
of $x$, so $\rk(xy)=r$ if and only if the first case occurs for all~$k$, that
is, if and only if $\bt(x)\sim\tp(y)$. In that case the blocks of $xy$ inside
the top row are exactly those of $x$, so $\tp(xy)=\tp(x)$. The $r$ strands end
at $r$ distinct bottom through points of $y$, that is, at all of them, so the
blocks of $xy$ inside the bottom row are those of $y$ and $\bt(xy)=\bt(y)$.
Following $a_k$: it reaches the $(k\,g(x))$-th smallest point of
$F(\bt(x))$, which is joined to the $(k\,g(x)\,p)$-th smallest point of
$F(\tp(y))$, where $p=p_{\bt(x),\tp(y)}$, and this point is joined in $y$ to
the $(k\,g(x)\,p\,g(y))$-th smallest bottom through point of~$y$.

(c) In $\Gamma(\lambda,\lambda)$ each arc of $\lambda$ is a double edge and each
point of $F(\lambda)$ is isolated, which gives $\lambda\sim\lambda$ and
$p_{\lambda\lambda}=1$. The graphs $\Gamma(\lambda,i)$ and $\Gamma(i,\lambda)$
coincide. If $\lambda\sim i$, the $r$ paths starting in $F(\lambda)$ end at $r$
distinct points of $F(i)$, hence at all of them; so the path starting at a
point of $F(i)$ is one of these paths, read backwards. Hence $i\sim\lambda$ and
$p_{i\lambda}=p_{\lambda i}^{-1}$.

(d) Each point of $F(\lambda)=F(i)$ lies on no arc of $\lambda$ or of $i$, so it
is a path of length~$0$.

Finally, $J_r$ is a $\mathcal J$-class \cite[Proposition~2.1]{emrt}, and
$\rk(xy)\le r$ for $x,y\in J_r$. So by (b) the product of $x,y\in J_r$ in
$J_r^0$ is $xy$ if $\bt(x)\sim\tp(y)$ and $0$ otherwise, and by (a) and~(b)
the stated map is an isomorphism. By~(c), $P$ has no zero row or column.
\end{proof}

The rows of $P$ are indexed by bottom halves ($\mathcal L$-classes) and the
columns by top halves ($\mathcal R$-classes). Proposition~\ref{prop:rees} is
the analogue for $\B_n$ and $\PB_n$ of \cite[Proposition~8.3]{abchk} for the
partition monoid.

For $\lambda\in H$, let $e_\lambda\in J_r$ be the element with coordinates
$(\lambda,1,\lambda)$, that is, the diagram whose top and bottom halves are
both $\lambda$ and whose through blocks are the pairs $\{a,a'\}$ with $a\in
F(\lambda)$. Since $p_{\lambda\lambda}=1$, Proposition~\ref{prop:rees}(b) gives
$(\lambda,g,\lambda)(\lambda,h,\lambda)=(\lambda,gh,\lambda)$ for $g,h\in
S_r$. So $e_\lambda$ is an idempotent, and its $\mathcal H$-class
$\{(\lambda,g,\lambda):g\in S_r\}$ is a group isomorphic to~$S_r$.

\section{Explicit complete mappings}\label{sec:cm}
Throughout this section $H$ is a non-empty finite set, $r\ge0$, and $P$ is an
$H\times H$ matrix over $S_r\cup\{0\}$ with $p_{\lambda\lambda}=1$ for all
$\lambda$. We write $\lambda\sim\mu$ when $p_{\lambda\mu}\ne0$. By
Proposition~\ref{prop:rees}, the principal factors of $\B_n$ and $\PB_n$ are
of this form. We fix a Latin square on $H$, that is, a map $L\colon H\times
H\to H$ such that $L(i,\cdot)$ and $L(\cdot,\lambda)$ are bijections for all
$i,\lambda$; for instance, identify $H$ with $\Z_N$ and put
$L(i,\lambda)=i+\lambda$.

\begin{lemma}\label{lem:A}
If $S_r$ has a complete mapping $\varphi$, then the map $f$ with $0f=0$ and
$(i,g,\lambda)f=(\lambda,\,g\varphi,\,L(i,\lambda))$ is a complete mapping of
$\M^0(S_r,H,H,P)$.
\end{lemma}

\begin{proof}
The map $(i,\lambda)\mapsto(\lambda,L(i,\lambda))$ is a bijection of $H\times
H$ and $\varphi$ is a bijection, so $f$ is a bijection. Since
$p_{\lambda\lambda}=1$, we get $(i,g,\lambda)\cdot(i,g,\lambda)f=(i,\,g\cdot
g\varphi,\,L(i,\lambda))\ne0$, and this is a bijection because
$(i,\lambda)\mapsto(i,L(i,\lambda))$ and $g\mapsto g\cdot g\varphi$ are
bijections.
\end{proof}

Lemma~\ref{lem:A} uses only the diagonal of $P$, as in
\cite[Proposition~6.5]{abchk}; the existence statement is a case of~(R6).

Now let $r\in\{2,3\}$. Let $s\colon S_r\to\Z_2=\{0,1\}$ be the sign
homomorphism written additively, with $s(g)=0$ exactly for even $g$; its kernel
is the alternating group $A_r$. We write $\Mt$ for the Rees $0$-matrix
semigroup $\M^0(\Z_2,H,H,s(P))$, with $\Z_2$ written additively: its non-zero
elements are the triples $(i,a,\lambda)\in H\times\Z_2\times H$, and
\[
 (i,a,\lambda)(j,b,\mu)=(i,\,a+s(p_{\lambda j})+b,\,\mu)\ \text{ if
 }\lambda\sim j,\qquad (i,a,\lambda)(j,b,\mu)=0\ \text{ otherwise.}
\]
Note that $s(p_{\lambda\lambda})=0$.

\begin{lemma}[lifting]\label{lem:lift}
Let $r\in\{2,3\}$, and suppose that $\Mt$ has a complete mapping. Then
$\M^0(S_r,H,H,P)$ also has a complete mapping.
\end{lemma}

\begin{proof}
The group $A_r$ has odd order ($1$ or $3$), so $h\mapsto h^2$ is a bijection of
$A_r$ and the identity map is a complete mapping of~$A_r$. Since
$S_r/A_r\cong\Z_2$ via~$s$, (R5) applies with $N=A_r$. Explicitly, following
the proof of \cite[Theorem~5.2]{abchk}: put $g_0=1$ and $g_1=(1\,2)$, and write
each $g\in S_r$ as $g=g_uh$ with $u=s(g)$ and $h\in A_r$. If $\bar f$ is a
complete mapping of $\Mt$ and $(i,u,\lambda)\bar f=(j,v,\mu)$, put
$(i,g_uh,\lambda)f=(j,\,p^{-1}hp\,g_v,\,\mu)$ with $p=p_{\lambda j}$, and $0f=0$.
Then $f$ is a bijection, and
$(i,g_uh,\lambda)\cdot(i,g_uh,\lambda)f=(i,\,g_uh^2p\,g_v,\,\mu)$; for fixed
$(i,u,\lambda)$ this runs bijectively through the coset of $A_r$ determined by
the product $(i,u,\lambda)\cdot(i,u,\lambda)\bar f$ in~$\Mt$.
\end{proof}

\begin{lemma}\label{lem:B}
Let $r\in\{2,3\}$. Let $\delta$ be a permutation of $H$ and $t\colon H\to\Z_2$
a map with $t(\mu\delta)=t(\mu)+1$ for all $\mu\in H$. Define $\bar f$ on $\Mt$
by $0\bar f=0$ and, writing $\mu=L(i,\lambda)$,
\[
 (i,\,1+t(\mu),\,\lambda)\bar f=(\lambda,\,1,\,\mu),\qquad
 (i,\,t(\mu),\,\lambda)\bar f=(\lambda,\,0,\,\mu\delta).
\]
Then $\bar f$ is a complete mapping of $\Mt$. Such $\delta$ and $t$ exist when
$N=|H|$ is even, so in that case $\M^0(S_r,H,H,P)$ has a complete mapping.
\end{lemma}

\begin{proof}
The element $(\lambda,1,\nu)$ is the image of exactly one element, coming from
the unique $i$ with $L(i,\lambda)=\nu$; the element $(\lambda,0,\nu)$ is the
image of exactly one element, coming from the unique $i$ with
$L(i,\lambda)\delta=\nu$. So $\bar f$ is a bijection. Since
$s(p_{\lambda\lambda})=0$, the products are
\[
 (i,1+t(\mu),\lambda)(\lambda,1,\mu)=(i,\,t(\mu),\,\mu),\qquad
 (i,t(\mu),\lambda)(\lambda,0,\mu\delta)=(i,\,t(\mu),\,\mu\delta)=
 (i,\,t(\mu\delta)+1,\,\mu\delta).
\]
For fixed $i$, $\mu=L(i,\lambda)$ runs through $H$ as $\lambda$ does, so the
products are the elements $(i,t(\nu),\nu)$ and $(i,t(\nu)+1,\nu)$ for $\nu\in
H$, each exactly once. If $N$ is even, let $\delta$ be a fixed-point-free
involution of $H$ and let $t$ take the values $0$ and $1$ on the two points of
each $2$-cycle. The last statement follows from Lemma~\ref{lem:lift}.
\end{proof}

Since the diagonal of $P$ is a one-transversal of its pattern, the existence
statement of Lemma~\ref{lem:B} is also a case of~(R7); the lemma gives an
explicit formula.

\begin{lemma}\label{lem:C}
Let $r\in\{2,3\}$. Let $\rho$ be a permutation of $H$ with $\lambda\sim\lambda\rho$
for all $\lambda\in H$ such that, for every cycle
$(\lambda_0,\lambda_1,\dots,\lambda_{m-1})$ of $\rho$, with indices modulo~$m$,
\begin{equation}\tag{$\ast$}\label{eq:star}
 m+\sum_{t=0}^{m-1}s(p_{\lambda_t\lambda_{t+1}})\equiv0\pmod2.
\end{equation}
Then there is a map $c\colon H\to\Z_2$ with
$c(\lambda\rho)=c(\lambda)+1+s(p_{\lambda,\lambda\rho})$ for all $\lambda$. For
such $c$, the map $\bar f$ with $0\bar f=0$ and
\[
 (i,\,c(\lambda),\,\lambda)\bar f=(\lambda,\,1+c(\lambda),\,L(i,\lambda)),\qquad
 (i,\,1+c(\lambda),\,\lambda)\bar f=(\lambda\rho,\,c(\lambda\rho),\,
 L(i,\lambda\rho))
\]
is a complete mapping of $\Mt$, and hence $\M^0(S_r,H,H,P)$ has a complete
mapping.
\end{lemma}

\begin{proof}
On each cycle of $\rho$, choose $c(\lambda_0)$ arbitrarily and define
$c(\lambda_{t+1})=c(\lambda_t)+1+s(p_{\lambda_t\lambda_{t+1}})$ for
$t=0,\dots,m-2$; condition \eqref{eq:star} says exactly that the required
relation also holds for $t=m-1$. The images under the first rule are the
elements $(\lambda,1+c(\lambda),\mu)$, and those under the second rule the
elements $(\kappa,c(\kappa),\mu)$ with $\kappa=\lambda\rho$; in both cases
$\mu$ runs through $H$ as $i$ does. These two sets are complementary, so $\bar
f$ is a bijection. The products are
\[
 (i,\,c(\lambda)+0+1+c(\lambda),\,L(i,\lambda))=(i,1,L(i,\lambda))
\]
from the first rule, which is non-zero since $\lambda\sim\lambda$, and
\[
 (i,\,1+c(\lambda)+s(p_{\lambda,\lambda\rho})+c(\lambda\rho),\,
 L(i,\lambda\rho))=(i,0,L(i,\lambda\rho))
\]
from the second, which is non-zero since $\lambda\sim\lambda\rho$. For fixed
$i$ they fill $\{i\}\times\{1\}\times H$ and $\{i\}\times\{0\}\times H$
bijectively. The last statement follows from Lemma~\ref{lem:lift}.
\end{proof}

\begin{remark}\label{rem:C}
(1) A fixed point of $\rho$ would be a cycle with $m=1$ and sum $0$, so $\rho$
has no fixed points. (2) If $\lambda\sim\mu$ implies $\mu\sim\lambda$ and
$p_{\mu\lambda}=p_{\lambda\mu}^{-1}$, as in Proposition~\ref{prop:rees}(c),
then every transposition $(a\ b)$ with $a\sim b$ satisfies \eqref{eq:star},
because $s(p_{ab})=s(p_{ba})$. (3) A $3$-cycle $(a\ b\ c)$ satisfies
\eqref{eq:star} if and only if $s(p_{ab})+s(p_{bc})+s(p_{ca})=1$, that is,
if and only if $p_{ab}p_{bc}p_{ca}$ is odd. (4) Lemma~\ref{lem:C} is the
special case of the residual-routing criterion \cite[Definition~7.13 and
Theorem~7.14]{abchk} in which the normal subgroup is $A_r$, the routing is
$(i,\lambda)\mapsto(\lambda,L(i,\lambda))$ and the residual permutation is
$(i,\lambda)\mapsto(i,\lambda\rho)$: the cycles of the residual permutation are
the sets $\{i\}\times C$ for the cycles $C$ of~$\rho$, and the cycle condition
of \cite{abchk} becomes~\eqref{eq:star}. We give the direct proof because it
yields the explicit formula that is formalised in Section~\ref{sec:lean}.
\end{remark}

\section{The Brauer monoid}\label{sec:odd}
In this section $D=\B_n$, $r\in\{2,3\}$, $r<n$ and $n\equiv r\pmod 2$, and
$H=H_r(\B_n)$. For an $r$-subset $F$ of $[n]$, the \emph{fibre} over $F$ is
$H_F=\{h\in H:F(h)=F\}$. It has $m=(n-r-1)!!$ elements, an odd number, and by
Proposition~\ref{prop:rees}(d) any two of its elements are compatible.

\begin{lemma}\label{lem:parity}
$N=|H|$ is odd if and only if $r=2$ and $n\equiv2\pmod4$, or $r=3$ and
$n\equiv3\pmod4$. In these cases $n-r\ge4$ and $\binom nr$ is odd.
\end{lemma}

\begin{proof}
Since $m$ is odd, $N=\binom nr m$ is odd if and only if $\binom nr$ is odd. For
even $n$, $\binom n2=\frac n2(n-1)$ is odd if and only if $n\equiv2\pmod4$. For
odd $n$, the $2$-adic valuation of $n(n-1)(n-2)$ is that of $n-1$, so
$\binom n3=n(n-1)(n-2)/6$ is odd if and only if $n-1\equiv2\pmod4$, that is,
$n\equiv3\pmod4$. Since $r<n$, this gives $n\ge6$ for $r=2$ and $n\ge7$ for
$r=3$, so $n-r\ge4$.
\end{proof}

\begin{lemma}[the triangle]\label{lem:triangle}
Suppose that $n-r\ge2$. Let $E=\varnothing$ if $r=2$ and $E=\{5\}$ if $r=3$, and
let $R_0$ be the perfect matching $\{r+3,r+4\},\{r+5,r+6\},\dots$ of
$\{r+3,\dots,n\}$. Put
\begin{align*}
 I_0&=(\{3,4\}\cup E,\ \{\{1,2\}\}\cup R_0),\\
 I_1&=(\{2,4\}\cup E,\ \{\{1,3\}\}\cup R_0),\\
 I_2&=(\{2,3\}\cup E,\ \{\{1,4\}\}\cup R_0).
\end{align*}
Then $I_0\sim I_1$, $I_1\sim I_2$ and $I_2\sim I_0$, with
$p_{I_0I_1}=p_{I_1I_2}=1$ and $p_{I_2I_0}=(1\,2)$. In particular
$p_{I_0I_1}p_{I_1I_2}p_{I_2I_0}$ is odd.
\end{lemma}

\begin{proof}
The arcs of $R_0$ are common to the three half-diagrams and give double edges,
and the point $5$ (if $r=3$) is free in all three and is the largest free
point, so it is joined to itself. Only the points $1,\dots,4$ matter. In
$\Gamma(I_0,I_1)$ the path from $3$ is $3,1,2$ (arcs $\{1,3\}$ of $I_1$ and
$\{1,2\}$ of $I_0$) and ends at $2\in F(I_1)$, and $4$ is fixed; the order is
preserved. In $\Gamma(I_1,I_2)$, $2$ is fixed and the path from $4$ is $4,1,3$;
the order is preserved. In $\Gamma(I_2,I_0)$ the path from $2$ is $2,1,4$ and
$3$ is fixed, so the smaller free point $2$ of $I_2$ is joined to the larger
free point $4$ of $I_0$, and $p_{I_2I_0}=(1\,2)$.
\end{proof}

These are the three half-diagrams used for $\mathcal P_n$ in the proof of
\cite[Theorem~8.5]{abchk}, where the marked singletons become free points,
completed to perfect matchings by~$R_0$. With the idempotents $e_\lambda$
defined after Proposition~\ref{prop:rees}, Lemma~\ref{lem:triangle} says that
$e_{I_0}e_{I_1}e_{I_2}e_{I_0}=(I_0,\,p_{I_0I_1}p_{I_1I_2}p_{I_2I_0},\,I_0)$
is an odd element of the group $\mathcal H$-class of~$e_{I_0}$.

\begin{lemma}[joining two fibres]\label{lem:join}
Let $F\ne F'$ be $r$-subsets of $[n]$ with $F\setminus F'=\{x_1,\dots,x_k\}$
and $F'\setminus F=\{y_1,\dots,y_k\}$, and suppose that $r+2k\le n$. Then
there are $v\in H_F$ and $w\in H_{F'}$ with $v\sim w$.
\end{lemma}

\begin{proof}
Choose distinct points $z_1,\dots,z_k\notin F\cup F'$, which is possible since
$|F\cup F'|+k=r+2k\le n$, and a perfect matching $R$ of the remaining
$n-r-2k$ points, an even number. Put
\[
 v=(F,\ \{\{z_t,y_t\}:t\le k\}\cup R),\qquad w=(F',\ \{\{x_t,z_t\}:t\le k\}\cup
 R).
\]
In $\Gamma(v,w)$ the path from $x_t$ is $x_t,z_t,y_t$ and ends in $F'$, the
points of $F\cap F'$ are paths of length~$0$, and the arcs of $R$ give double
edges. So $v\sim w$.
\end{proof}

\begin{lemma}[matching]\label{lem:matching}
Suppose that $n-r\ge4$ and $\binom nr$ is odd, and let $X=\{I_0,I_1,I_2\}$ be
as in Lemma~\ref{lem:triangle}. Then $H\setminus X$ can be partitioned into
pairs $\{v,w\}$ with $v\sim w$.
\end{lemma}

\begin{proof}
The elements of $X$ lie in the three distinct fibres over $\{3,4\}\cup E$,
$\{2,4\}\cup E$ and $\{2,3\}\cup E$, one in each. The other $\binom nr-3$
fibres, an even number, can be paired so that each pair $(F,F')$ satisfies
$r+2|F\setminus F'|\le n$. Indeed, if $r=2$ then $|F\setminus F'|\le2$ and
$n\ge6$, and if $r=3$ and $n\ge9$ then $|F\setminus F'|\le3$; in these cases
any pairing works. If $(n,r)=(7,3)$ the condition is $F\cap F'\ne\varnothing$,
and the following $16$ pairs of $3$-subsets of $[7]$ work: $\{a,b,6\}$ with
$\{a,b,7\}$ for the ten $2$-subsets $\{a,b\}$ of $[5]$; $\{1,6,7\}$ with
$\{1,2,3\}$, $\{2,6,7\}$ with $\{1,2,4\}$, $\{3,6,7\}$ with $\{1,3,4\}$,
$\{4,6,7\}$ with $\{2,3,4\}$ and $\{5,6,7\}$ with $\{1,2,5\}$; and $\{1,3,5\}$
with $\{1,4,5\}$. These are exactly the $32$ subsets other than $\{3,4,5\}$,
$\{2,4,5\}$ and $\{2,3,5\}$, and the two sets in each pair meet. For each
pair of fibres, Lemma~\ref{lem:join} gives compatible $v\in H_F$ and $w\in
H_{F'}$, and we take $\{v,w\}$ as a pair. Now exactly one element of each fibre
has been used, either in $X$ or in a pair. The remaining $m-1$ elements of each
fibre, an even number, are pairwise compatible, and we pair them arbitrarily.
\end{proof}

\begin{proof}[Proof of Theorem~\ref{thm:main}(a)]
Let $r<n$ be a rank that occurs in $\B_n$. By Proposition~\ref{prop:rees},
$J_r^0\cong\M^0(S_r,H,H,P)$ with $p_{\lambda\lambda}=1$ for all $\lambda$. If
$r\notin\{2,3\}$, then $S_r$ has a complete mapping by (R1), and
Lemma~\ref{lem:A} applies. If $r\in\{2,3\}$ and $N$ is even,
Lemma~\ref{lem:B} applies. If $r\in\{2,3\}$ and $N$ is odd, then $n-r\ge4$ and
$\binom nr$ is odd by Lemma~\ref{lem:parity}. Let $\rho$ be the permutation of
$H$ that maps $I_0\mapsto I_1\mapsto I_2\mapsto I_0$ and interchanges the two
elements of each pair given by Lemma~\ref{lem:matching}. Then
$\lambda\sim\lambda\rho$ for all $\lambda$. The $3$-cycle satisfies
\eqref{eq:star} by Lemma~\ref{lem:triangle} and Remark~\ref{rem:C}(3), and the
transpositions satisfy it by Proposition~\ref{prop:rees}(c) and
Remark~\ref{rem:C}(2). So Lemma~\ref{lem:C} gives a complete mapping of
$J_r^0$.
\end{proof}

\section{The partial Brauer monoid}\label{sec:pb}

\begin{lemma}\label{lem:pbparity}
The number $a(k)$ of involutions of a $k$-set is even for $k\ge2$. Hence, if
$n\ge4$, $r\in\{2,3\}$ and $r<n$, then $N=\binom nr a(n-r)$ is even.
\end{lemma}

\begin{proof}
We have $a(2)=2$, $a(3)=4$ and $a(k)=a(k-1)+(k-1)a(k-2)$, so $a(k)$ is even for
$k\ge2$ by induction. If $n-r\ge2$ this gives the claim; otherwise $n-r=1$, so
$(n,r)=(4,3)$ and $N=4$.
\end{proof}

\begin{proposition}\label{prop:pb3}
The principal factor of rank~$2$ of $\PB_3$ is isomorphic to the Brandt
semigroup $B(S_2,3)=\M^0(S_2,[3],[3],\Delta_3)$ and has no complete mapping.
\end{proposition}

\begin{proof}
The half-diagrams of rank $2$ are the three pairs $(F,\varnothing)$ with
$|F|=2$, the third point being a singleton. If $F(\lambda)\ne F(i)$, a point of
$F(\lambda)\setminus F(i)$ is a singleton of $i$, hence an isolated vertex of
$\Gamma(\lambda,i)$ outside $F(i)$, so $\lambda\not\sim i$. With
Proposition~\ref{prop:rees}(c), the sandwich matrix is the identity matrix
$\Delta_3$ over $S_2$. It is obtained from the normalised $3\times3$ matrix all
of whose entries are the identity, which has odd order, by replacing the
off-diagonal entries by~$0$. Since $S_2\cong C_2$ has a non-trivial cyclic
Sylow $2$-subgroup and $3$ is odd, (R8) shows that there is no complete
mapping.
\end{proof}

The same conclusion follows from \cite[Theorem~10.1]{abchk}, since $B(S_2,3)$
is an inverse semigroup whose non-zero $\mathcal J$-class has three $\mathcal
L$-classes and maximal subgroup~$S_2$. An exhaustive computer search confirms
it (Section~\ref{sec:comp}).

\begin{proof}[Proof of Theorem~\ref{thm:main}(b)]
Let $r<n$. If $r\notin\{2,3\}$, Lemma~\ref{lem:A} applies as in the proof of
part~(a). If $r\in\{2,3\}$ and $n\ge4$, then $N$ is even by
Lemma~\ref{lem:pbparity}, and Lemma~\ref{lem:B} applies. If $r\in\{2,3\}$ and
$n\le3$, then $(n,r)=(3,2)$, and Proposition~\ref{prop:pb3} applies.
\end{proof}

\section{The whole monoids}\label{sec:whole}

\begin{proof}[Proof of Theorem~\ref{thm:main}(c)]
The group of units of $\B_n$ and of $\PB_n$ is $J_n\cong S_n$, so (i) and (ii)
each imply (iii) by (R3). Statements (iii) and (iv) are equivalent by (R1).
Suppose that (iii) holds, and let $\varphi$ be a complete mapping of $S_n$. The
principal factor of the group of units is $S_n$ with a zero adjoined, and
$\varphi$ extended by $0\mapsto0$ is a complete mapping of it. Every proper
principal factor of $\B_n$ has a complete mapping by part~(a). Every proper
principal factor of $\PB_n$ has one by part~(b), since $n\ne3$. So (i) and
(ii) hold by (R4).
\end{proof}

\section{Computations}\label{sec:comp}
No statement of this paper depends on a computation: each has a written proof
above. The computations are independent checks of these proofs and of our
conventions. The programs (in Python), certificates and logs are in the
directory \texttt{certificates/brauer-complete} of version~1.4.0 of the
repository~\cite{repo}.

\paragraph{Two implementations.}
The first implementation multiplies diagrams with a union--find structure,
computes the Rees coordinates of Proposition~\ref{prop:rees}, and builds
complete mappings exactly as in the proofs (Lemmas~\ref{lem:A}--\ref{lem:C}
and~\ref{lem:triangle}--\ref{lem:matching}, with $L(i,\lambda)=i+\lambda$ after
indexing $H$ by~$\Z_N$). The second implementation shares no code with the
first: it has its own parser, multiplies diagrams by following alternating
paths through the middle row, enumerates the elements itself, and never uses
Rees coordinates. A certificate is a text file with one line
$x\mapsto x\alpha$ for every element $x$ of a principal factor (without the
zero, which is fixed) or of a whole monoid. For a factor $J_r^0$, the second
implementation checks that the listed $x$ are valid diagrams of rank~$r$,
pairwise distinct, and $N^2r!=|J_r|$ in number, so that they are exactly the
elements of $J_r$; that the values $x\alpha$ are pairwise distinct and of
rank~$r$; and that the products $x\cdot x\alpha$ have rank~$r$ (that is, are
non-zero in $J_r^0$) and are pairwise distinct. For a whole monoid it checks
that the listed $x$ are all the elements, that $\alpha$ is a bijection, and
that $x\mapsto x\cdot x\alpha$ is a bijection. Table~\ref{tab:certs} lists the
certificates; all passed.

\begin{table}[ht]\centering\small
\begin{tabular}{@{}l r >{\raggedright\arraybackslash}p{0.56\textwidth}@{}}
\toprule
Certificate for & elements & construction\\\midrule
$J_2^0(\B_6)$, $N=45$ & $4050$ & Lemmas~\ref{lem:C}, \ref{lem:triangle},
 \ref{lem:matching} and~\ref{lem:lift}; the first case with $N$ odd\\
$J_3^0(\B_7)$, $N=105$ & $66150$ & the same; the first case with $N$ odd and
 $r=3$\\
$\B_6$ & $10395$ & ranks $0$, $2$, $4$, $6$; rank~$2$ as above\\
$\B_1$, $\B_4$, $\B_5$ & $1$, $105$, $945$ & Lemmas~\ref{lem:A},
 \ref{lem:lift} and~\ref{lem:B}\\
$\PB_6$ & $140152$ & ranks $0$ to $6$; ranks $2$ and $3$ by Lemma~\ref{lem:B}\\
$\PB_1$, $\PB_4$, $\PB_5$ & $2$, $764$, $9496$ & Lemmas~\ref{lem:A},
 \ref{lem:lift} and~\ref{lem:B}\\\addlinespace
$J_2^0(\B_6)$, $J_3^0(\B_7)$ & $4050$, $66150$ & second versions: a
 triangle with sign bits $(0,1,0)$ found by search and a maximum matching
 computed with the NetworkX library in place of
 Lemmas~\ref{lem:triangle} and~\ref{lem:matching}\\
\bottomrule
\end{tabular}
\caption{Certificates. The number of elements excludes the zero of a
principal factor. Each certificate was built by the first implementation and
checked by the second.}
\label{tab:certs}
\end{table}

For the whole monoids, the factors of rank at least $4$ and the groups of units
need complete mappings of $S_4$, $S_5$ and $S_6$. Those of $S_4$ and $S_5$ were
found with the SAT solver CaDiCaL, and one of $S_6$ by an exact-cover search
(Knuth's Algorithm~X) in about seven minutes, after CaDiCaL had not finished
in twenty minutes. They are checked implicitly by the whole-monoid
certificates. For the theorems themselves, (R1) supplies these complete
mappings.

\paragraph{Further checks.}
\begin{itemize}
\item \emph{Rees coordinates.} The first implementation confirms that the
 $\mathcal J$-classes are the rank classes for $\B_n$ with $n\le5$ and $\PB_n$
 with $n\le4$, and that Proposition~\ref{prop:rees}(a)--(c) and the product
 rule~(b) hold for all pairs in $J_r$ for these monoids. For $\B_6$ the product
 rule was checked for all pairs in ranks $0$ and $6$ and for $400\,000$ random
 pairs in each of ranks $2$ and~$4$; for $\PB_5$, for all pairs in the ranks
 with $|J_r|^2\le4\cdot10^6$ and for $400\,000$ random pairs in each other
 rank. The values of $N$ agree with the formulas of Section~\ref{sec:rees}.
\item \emph{The odd case.} For $(n,r)=(6,2)$, $(7,3)$, $(10,2)$, $(11,3)$ and
 $(14,2)$, that is, for $N=45$, $105$, $4725$, $17325$ and $945945$, the first
 implementation confirms the sandwich permutations of
 Lemma~\ref{lem:triangle} and that the construction of
 Lemma~\ref{lem:matching} yields a partition of $H\setminus X$ into
 compatible pairs. The second implementation confirms, without Rees
 coordinates, that for $n\in\{6,7,10,11\}$ the product
 $e_{I_0}e_{I_1}e_{I_2}e_{I_0}$, with $e_\lambda$ as defined after
 Proposition~\ref{prop:rees}, is the element of the group $\mathcal H$-class
 of $e_{I_0}$ that interchanges the through strands at the points $3$ and~$4$.
\item \emph{Parities.} The parity pattern of Table~\ref{tab:cases} holds for
 all $n\le200$ in both families.
\item \emph{Non-existence.} An exhaustive search, using the second
 implementation, confirms that $S_2$ and $S_3$ have no complete mapping, that
 the units of $\B_2$, $\B_3$, $\PB_2$ and $\PB_3$ are the permutation diagrams,
 and that the $19$-element rank-$2$ principal factor of $\PB_3$ has no complete
 mapping (a depth-first search over all bijections, with $67\,911\,890$
 nodes). As positive controls, the same search finds complete mappings of
 three other small principal factors.
\end{itemize}
One script reruns all these checks and regenerates the certificates; on its
last run (26 September 2026, 07:42--07:44 UTC) every certificate was
reproduced with the same SHA-256 hash, and the hashes are recorded in the
repository.

\section{Formal verification}\label{sec:lean}
The Lean file \texttt{lean/Research/BrauerCMCore.lean} of version~1.4.0 of
the repository~\cite{repo} (324 lines) formalises Lemmas~\ref{lem:B}
and~\ref{lem:C} at the level of the semigroup $\Mt$, for an arbitrary finite
index set. It uses Lean~4 (version 4.33.1) and Mathlib~\cite{mathlib} (commit
\texttt{0df444a3}). The declarations are in the namespace
\texttt{BrauerCMCore}. The assertion of Lemma~\ref{lem:B} that $\bar f$ is a
complete mapping of $\Mt$ is the theorem \texttt{BrauerCMCore.lemmaB}, and the
corresponding assertion of Lemma~\ref{lem:C} is the theorem
\texttt{BrauerCMCore.lemmaC}. The following listing is an excerpt of this
file, not the complete file. It shows the definitions of \texttt{Elt},
\texttt{mul} and \texttt{IsCompleteMapping} and the statements of
\texttt{lemmaC} and \texttt{lemmaB}, each inside its section together with the
\texttt{variable} declaration of that section. Proofs, auxiliary lemmas,
comments and the definitions of \texttt{fC} and \texttt{fB} are omitted, and
some long lines are broken differently.

\vspace{12pt}
\noindent\begin{minipage}{\linewidth}
\begin{alltt}\small
abbrev Elt (H : Type*) := Option (H \(\times\) ZMod 2 \(\times\) H)
variable \{H : Type*\}
def mul (P : H \(\to\) H \(\to\) Option (ZMod 2)) : Elt H \(\to\) Elt H \(\to\) Elt H
  | none, _ => none
  | some _, none => none
  | some (i, a, l), some (j, b, m) => (P l j).map (fun p => (i, a + p + b, m))
def IsCompleteMapping \{S : Type*\} (op : S \(\to\) S \(\to\) S) (\(\alpha\) : S \(\to\) S) : Prop :=
  Function.Bijective \(\alpha\) \(\wedge\) Function.Bijective (fun x => op x (\(\alpha\) x))
\end{alltt}
\end{minipage}

\vspace{10pt}
\noindent\begin{minipage}{\linewidth}
\begin{alltt}\small
section LemmaC
variable (P : H \(\to\) H \(\to\) Option (ZMod 2)) (L : H \(\to\) H \(\to\) H) (\(\rho\) : H \(\simeq\) H)
  (c s : H \(\to\) ZMod 2)
theorem lemmaC [Finite H] (hdiag : \(\forall\) l, P l l = some 0)
    (h\(\rho\) : \(\forall\) l, P l (\(\rho\) l) = some (s l)) (hc : \(\forall\) l, c (\(\rho\) l) = c l + 1 + s l)
    (hrow : \(\forall\) i, Function.Injective (L i))
    (hcol : \(\forall\) l, Function.Injective (fun i => L i l)) :
    IsCompleteMapping (mul P) (fC L \(\rho\) c)
end LemmaC
\end{alltt}
\end{minipage}

\vspace{10pt}
\noindent\begin{minipage}{\linewidth}
\begin{alltt}\small
section LemmaB
variable (P : H \(\to\) H \(\to\) Option (ZMod 2)) (L : H \(\to\) H \(\to\) H) (\(\delta\) : H \(\simeq\) H)
  (t : H \(\to\) ZMod 2)
theorem lemmaB [Finite H] (hdiag : \(\forall\) l, P l l = some 0)
    (h\(\delta\) : \(\forall\) \(\mu\), t (\(\delta\) \(\mu\)) = t \(\mu\) + 1)
    (hrow : \(\forall\) i, Function.Injective (L i))
    (hcol : \(\forall\) l, Function.Injective (fun i => L i l)) :
    IsCompleteMapping (mul P) (fB L \(\delta\) t)
end LemmaB
\end{alltt}
\end{minipage}

\vspace{14pt}
\noindent Here \texttt{none} is the zero of $\Mt$, \texttt{some (i, a, l)} is the
element $(i,a,\lambda)$, and \texttt{P l j = none} means $\lambda\not\sim j$.
The maps \texttt{fC} and \texttt{fB} send \texttt{none} to \texttt{none} and
are the maps of Lemmas~\ref{lem:C} and~\ref{lem:B}: \texttt{fC} sends
\texttt{some (i, a, l)} to \texttt{some (l, 1~+~c~l, L~i~l)} if
\texttt{a~=~c~l}, and to \texttt{some ($\rho$~l, c~($\rho$~l),
L~i~($\rho$~l))} otherwise; \texttt{fB} sends it to \texttt{some (l, 1, L~i~l)} if
\texttt{a~=~1~+~t~(L~i~l)}, and to \texttt{some (l, 0, $\delta$~(L~i~l))}
otherwise. The predicate \texttt{IsCompleteMapping} is the definition of a
complete mapping of a magma; associativity is not used.

The Lean statements differ from Lemmas~\ref{lem:B} and~\ref{lem:C} as
follows. The matrix is over $\Z_2$ directly, with no reference to $S_r$ or to
diagrams; only its diagonal entries (\texttt{hdiag}) and, for
\texttt{lemmaC}, its entries at $(\lambda,\lambda\rho)$ (\texttt{h$\rho$}) are
constrained. In \texttt{lemmaC} the map $c$ is a hypothesis; that such a $c$
exists under condition~\eqref{eq:star} is the first step of the proof of
Lemma~\ref{lem:C} and is not formalised. The Latin square is given by the
hypotheses \texttt{hrow} and \texttt{hcol}.

The following are not formalised: Proposition~\ref{prop:rees} (diagram monoids
are not in Mathlib); the lifting Lemma~\ref{lem:lift}, that is, (R5);
Lemma~\ref{lem:A}; Lemmas~\ref{lem:parity}--\ref{lem:matching},
\ref{lem:pbparity} and Proposition~\ref{prop:pb3}; the results (R1)--(R8)
of~\cite{abchk} and the results of~\cite{emrt} that we use; and hence
Theorem~\ref{thm:main} itself. The Lean proof certifies the two constructions
of complete mappings over $\Z_2$ for every finite index set, in particular for
every $N$; it does not certify that the hypotheses hold for the principal
factors of $\B_n$ and $\PB_n$.

The file contains no \texttt{sorry} and no axiom declarations, and does not use
\texttt{native\_decide}; the ten uses of \texttt{decide} are closed statements
about $\Z_2$ with at most three variables. The module was built with
\texttt{lake build Research.BrauerCMCore} on 26 September 2026 (exit code~$0$).
The axiom audits (\texttt{\#print axioms}) of \texttt{lemmaC} and
\texttt{lemmaB} report only the standard axioms \texttt{propext},
\texttt{Classical.choice} and \texttt{Quot.sound}. The compiled module was
then checked with \texttt{leanchecker} (exit code~$0$), which replays its
declarations in Lean's kernel and trusts the imported Mathlib; this is not an
independent kernel. The code and logs are available at~\cite{repo}.

\section*{Tool and computational resource disclosure}
This work was carried out on 26 September 2026 (UTC) in a workflow directed by
the author, using AI coding agents: Anthropic Claude Code, with the model
Claude Opus~5.5. The author set the research programme, namely resolving open
problems with AI agents and formal verification, and directed the agents; the
agents selected this problem. Claude Code agents found the proofs, wrote the
programs and the Lean formalisation, carried out the literature searches
described in the introduction, and drafted this text. The programs, logs and
Lean code are available at~\cite{repo}. The author is not a professional
mathematician.
%% AUTHOR-ROLE-SENTENCE: to be completed after the author has read the paper
The correctness of the results rests on the written proofs, which use the
results of~\cite{abchk} and~\cite{emrt} listed in Section~\ref{sec:prelim}; the
Lean formalisation covers only the two constructions described in
Section~\ref{sec:lean}, and the computations of Section~\ref{sec:comp} are
additional checks. AI systems are not authors of this paper; the author takes
full responsibility for its content.

\begin{thebibliography}{9}
\bibitem{abchk}
J.~Ara\'ujo, W.~Bentz, P.~J. Cameron, K.~Hendrey and M.~Kinyon, Complete
mappings of semigroups, arXiv:2608.25092v1, 25 August 2026.
\bibitem{camblog}
P.~J. Cameron, Complete mappings of semigroups, \emph{Peter Cameron's Blog},
27 August 2026,
\url{https://cameroncounts.wordpress.com/2026/08/27/complete-mappings-of-semigroups/}.
\bibitem{repo}
H.~Dong, \emph{ai4math-results}, repository, version~1.4.0,
\url{https://github.com/Horace-Maxwell/ai4math-results}; archived at Zenodo,
concept DOI \url{https://doi.org/10.5281/zenodo.22970254}.
\bibitem{emrt}
J.~East, J.~D. Mitchell, N.~Ru\v{s}kuc and M.~Torpey, Congruence lattices of
finite diagram monoids, \emph{Adv.\ Math.} \textbf{333} (2018), 931--1003,
\url{https://doi.org/10.1016/j.aim.2018.05.016}.
\bibitem{kinyon}
M.~Kinyon, Complete mappings for semigroups, slides of a talk at AGTA,
University of Naples Federico II, 24 June 2025,
\url{https://www.advgrouptheory.com/agta2025/Talks/Kinyon.pdf}.
\bibitem{mathlib}
The mathlib Community, The Lean mathematical library, in \emph{Proceedings of
the 9th ACM SIGPLAN International Conference on Certified Programs and Proofs}
(CPP '20), ACM, 2020, 367--381,
\url{https://doi.org/10.1145/3372885.3373824}.
\bibitem{trurepo}
The Omega Institute, \texttt{trureturing} repository,
\url{https://github.com/the-omega-institute/trureturing}.
\end{thebibliography}
\end{document}
